\section{Detailed results}\label{sec:results}

\subsection{What is matched matters, how it is matched does not}\label{sec:results:factorial}

Figure~1(b) of the main paper shows the two factor effects of the 2$\times$2 design for each classifier. Replacing the boundary arm by the coverage arm raises the mean G-mean by 0.046 to 0.085 under the stable concept and by 0.050 to 0.086 under the optimal concept, depending on the classifier, and every 95\% bootstrap confidence interval excludes zero; the Wilcoxon signed-rank tests confirm the difference for all five classifiers under both concepts after Holm correction (Table~\ref{tab:factorial}; median differences 0.013 to 0.063). Replacing the stable by the optimal matching changes the mean G-mean by at most 0.003 in either arm; every confidence interval contains zero and lies within $\pm0.009$, and no Wilcoxon test approaches significance (smallest Holm-adjusted $p=0.342$ in the boundary arm and $0.507$ in the coverage arm). ROC-AUC gives the same picture, with the arm effect significant for four of the five classifiers.

The selections explain why the solution concept does not matter. \MUSBS{} and \MUSBO{} retain largely the same instances (median Jaccard similarity 0.91, against 0.03 expected by chance; Table~\ref{tab:overlap}), and over the 1375 training folds the stable selection is within a median 0.09\% (at most 4.7\%) of the optimal $\Wone$ distance to the minority class, which \MUSBO{} attains in every fold, as Proposition~3 of the main paper predicts. An exact transport solver that is independent of our implementation confirmed this optimum on the 275 folds of the first repetition (largest relative difference $6\times10^{-13}$). In the coverage arm the retained instances differ, because the member kept from each pair is chosen at random---\MUSCS{} and \MUSCO{} share as few instances as two runs of \MUSCS{} with different seeds (median Jaccard similarity 0.03)---but the retained distributions are nearly identical (median absolute differences of 0.005 in BP, 0.006 in CE and 0.003 in energy distance; these diagnostics are defined in Section~\ref{sec:results:mechanisms}). This happens even though the greedy matching is not optimal: on the first fold of every dataset, the greedy perfect matching of a round cost a median 6.5\% more than the optimal one on the same set (at most 27\% over 243 rounds; median 6.2\% and at most 46\% over 1096 rounds re-run on folds of the first repetition; Table~\ref{tab:a7a8}), far from the worst case of greedy matching \citep{reingold1981greedy}. Restricting the optimal matching to the candidate graph $G_{10}$ cost little: on the first fold of the 39 datasets with at most 1200 majority instances, its first-round optimum exceeded the complete-graph optimum by a median 0.03\% (at most 1.6\%).

Figure~1(c) of the main paper shows the same contrast at the level of datasets, as the G-mean difference to \RUS{} averaged over the five classifiers. The coverage arm lies within $\pm0.01$ of \RUS{} on 35 (\MUSCS) and 38 (\MUSCO) of the 55 datasets, above this region on 11 and below it on 9 and 6 datasets; its mean differences are $+0.0002$ and $+0.0015$. The boundary arm lies more than 0.01 below \RUS{} on 39 (\MUSBS) and 40 (\MUSBO) datasets and more than 0.01 above it on only 2.

\input{tables/stats_factorial_2x2}

\subsection{Comparison with the standard undersamplers}\label{sec:results:overall}

The Friedman and Iman--Davenport tests reject equal average ranks of the sixteen undersampling-family methods in all ten classifier--metric settings (largest $p$-values $8.1\times10^{-29}$ and $1.2\times10^{-32}$, respectively; Table~\ref{tab:ranks_friedman}). For G-mean, a coverage-arm method has the lowest average rank for every classifier (Figure~2b of the main paper): \MUSCO{} for CART (5.61), $k$-NN (4.94), SVM (4.11) and logistic regression (3.87), and \MUSCS{} for random forests (4.55). ClusterCentroids and \RUS{} follow closely, within the Nemenyi critical difference of 3.11 for every classifier (Figures~\ref{fig:cd_gmean} and~\ref{fig:cd_auc}). The boundary arm ranks between 6.70 and 9.71 (\MUSBS) and between 7.05 and 9.64 (\MUSBO), close to NearMiss-3 (6.98 to 9.51).

The pairwise tests sharpen this picture (Table~\ref{tab:wilcoxon_gmean} and Table~\ref{tab:wilcoxon_auc}). For G-mean, no matching undersampler is significantly better than \RUS, ClusterCentroids or the best non-matching method of any classifier after Holm correction. The coverage arm is never significantly worse either: its median differences to \RUS{} range from $+0.0005$ to $+0.0025$ (\MUSCS) and from $+0.0000$ to $+0.0021$ (\MUSCO). The boundary arm is significantly worse than all three comparators for all five classifiers; the median G-mean difference of \MUSBS{} to \RUS{} ranges from $-0.0190$ (random forests) to $-0.0583$ ($k$-NN). The Bayesian signed-rank tests (Figure~2a of the main paper; Tables~\ref{tab:bayes_gmean} and~\ref{tab:bayes_auc}) turn the absence of a significant difference into positive evidence of equivalence: with a region of practical equivalence of $\pm0.01$, \MUSCS{} and \MUSCO{} are practically equivalent to \RUS{} with posterior probability at least 0.95 for all five classifiers, in G-mean as well as in ROC-AUC. Against the best non-matching method they are practically equivalent for four classifiers; for CART, compared with the instance hardness threshold, the posterior is inconclusive (G-mean: most probably better; ROC-AUC: most probably worse). Against ClusterCentroids, practical equivalence in G-mean holds for SVM, logistic regression and random forests, whereas for CART and $k$-NN the posterior is inconclusive, with the remaining probability mostly on the side of the coverage arm; in ROC-AUC it holds for nine of the ten method--classifier pairs (Table~\ref{tab:a3_bayes_rope}). At narrower regions of practical equivalence the evidence for equivalence is weaker (Section~\ref{sec:results:robustness}). The boundary arm is worse than \RUS{} and than the best non-matching method with probability at least 0.95 for all five classifiers in G-mean.

ROC-AUC orders the methods differently. Methods that do not balance the classes lead for the stronger learners---Tomek links for SVM (average rank 5.67) and random forests (4.98), no resampling for logistic regression (4.86)---and the instance hardness threshold leads for CART (5.09). This is expected, because ROC-AUC does not depend on the decision threshold that balancing shifts, whereas every discarded instance is lost information for the ranking. Among the balancing methods, a coverage-arm method has the lowest ROC-AUC rank for every classifier. The coverage arm is significantly better than \RUS{} in ROC-AUC for some classifiers---\MUSCS{} for $k$-NN (median difference $+0.0022$, $p_\mathrm{H}=0.026$), SVM ($+0.0003$, $p_\mathrm{H}=0.041$) and logistic regression ($+0.0010$, $p_\mathrm{H}=0.006$), and \MUSCO{} for $k$-NN ($+0.0022$, $p_\mathrm{H}=0.030$) and random forests ($+0.0004$, $p_\mathrm{H}=0.008$)---but never better than ClusterCentroids, and it is worse than the best non-matching method for logistic regression (no resampling) and, for \MUSCS, random forests (Tomek links). These differences are statistically detectable but small: all median differences to \RUS{} are below 0.003, and those to the best non-matching method are at most 0.0045 in absolute value.

\input{tables/stats_ranks_friedman}
\input{tables/stats_wilcoxon_holm}

\subsection{Why the arms differ}\label{sec:results:mechanisms}

To characterise the selections we computed, for every training fold and undersampler, the \emph{boundary proximity} BP---the mean distance of the retained majority instances to their nearest minority instance, relative to the same quantity for all majority instances, so that BP $<1$ means that the retained instances lie closer to the minority class---and the \emph{coverage error} CE---the mean distance of every majority instance to its nearest retained instance, relative to the mean nearest-neighbour distance within the majority class. We also computed energy distances between the retained and the full majority class and the share of retained majority instances whose neighbourhoods are dominated by the minority class (Table~\ref{tab:diag_summary}). Because methods that keep more instances cover the majority class better mechanically, Figure~\ref{fig:mechanism}(a) compares only the balancing methods; the full map is given in Figure~\ref{fig:map_full}.

On this map the boundary arm lies next to NearMiss (median BP 0.67 and CE 2.47 for both \MUSBS{} and \MUSBO, against 0.58 to 0.82 and 2.62 to 3.03 for NearMiss and 1.00 and 1.81 for \RUS). It keeps majority instances closer to the minority class than \RUS{} in 54 of the 55 datasets but covers the majority class better in only 2, and it retains majority instances inside minority regions at 6.8 times their base rate (\RUS: 1.03; computed on the 40 datasets where such instances exist). The coverage arm lies between \RUS{} and ClusterCentroids (median BP 1.02, CE 1.56) and is closest to the full majority distribution among all balancing methods (median energy distance 0.107 for \MUSCS{} and 0.110 for \MUSCO, against 0.169 for ClusterCentroids and 0.170 for \RUS). The exact Wasserstein distance between the retained and the full majority class, computed on the folds of the first repetition, gives the same picture: it is a median 9\% (\MUSCS) and 7\% (\MUSCO) lower than for \RUS{} and lower on 50 and 48 of the 55 datasets, it is on par with \RUS{} for ClusterCentroids (median ratio 1.00, lower on 26 datasets), and it is 60 to 81\% higher for NearMiss and the boundary arm (Table~\ref{tab:a7a8}). The bound of Corollary~1 of the main paper held in every fold but exceeded the realised distance by a median factor of 2.6 for \MUSCS{} and 2.5 for \MUSCO{} (on 249 folds whose \MUSCO{} run took less than two minutes). The partial-round term contributed a median 35\% (\MUSCS) and 37\% (\MUSCO) of the bound, and the summed matching costs alone exceeded the realised distance in 233 of the 275 folds of \MUSCS, because chaining the rounds by the triangle inequality overcounts the transport.

Proximity on the map does not imply the same failure. Figure~\ref{fig:mechanism}(b) decomposes the change relative to \RUS{} into sensitivity and specificity, averaged over datasets and classifiers. NearMiss-1 and NearMiss-2 lose mainly specificity ($-0.114$ and $-0.184$, with sensitivity changes of $-0.043$ and $-0.004$), and NearMiss-3 loses both ($-0.056$ in sensitivity, $-0.052$ in specificity). The boundary arm loses sensitivity ($-0.111$ for both methods) and gains a little specificity ($+0.018$). Paired tests confirm that the routes differ: compared with NearMiss-1, \MUSBS{} has a higher specificity (median differences $+0.017$ to $+0.183$) and a lower sensitivity ($-0.011$ to $-0.040$) for all five classifiers (Wilcoxon, Holm-corrected over the ten tests, $p_\mathrm{H}\le0.006$). One explanation consistent with these routes is exclusivity: NearMiss lets many minority instances claim the same majority instances and leaves much of the majority region unrepresented, whereas the boundary arm gives every minority instance, including outliers and rare examples, its own majority partner, which may then claim the minority neighbourhood for the majority class. The experiments at the end of this subsection test this explanation. The coverage arm and ClusterCentroids make a much smaller trade in the same direction as the boundary arm (sensitivity $-0.009$ for \MUSCS, $-0.007$ for \MUSCO{} and $-0.014$ for ClusterCentroids; specificity $+0.011$ for both coverage methods and $+0.014$ for ClusterCentroids), consistent with their small advantages over \RUS{} in the precision-sensitive F1 and MCC (Section~\ref{sec:results:robustness}).

Which property of a selection matters for accuracy is harder to establish, because the diagnostics are confounded. Across the balancing methods, G-mean correlates within datasets with both coverage (median Spearman correlation with CE $-0.82$) and boundary proximity (with BP $+0.65$), but BP and CE are themselves strongly correlated (median $-0.80$), so these correlations cannot apportion the effect (Figure~\ref{fig:mechanism}c). Restricting the set to one method per arm and excluding NearMiss-3 (\RUS, NearMiss-1 and -2, ClusterCentroids, \MUSBS{} and \MUSCS) leaves this unchanged (median correlations $+0.71$ with BP, $-0.80$ with CE and $-0.83$ between BP and CE; G-mean correlations on 54 datasets because all six methods reach a G-mean of 1 on one dataset). Two contrasts are more informative, although both are weak evidence. Within the boundary-focused methods (NearMiss and the boundary arm), where the two diagnostics are only weakly related (median $-0.30$), G-mean tracks coverage (median $-0.70$, Holm-adjusted $p=2.2\times10^{-8}$) but not proximity ($+0.10$, $p_\mathrm{H}=0.32$); this contrast, however, rests on five methods, two of which select nearly the same instances and one of which keeps fewer instances than the target on 26 datasets. The $p$-values of the within-dataset correlations come from one-sample Wilcoxon signed-rank tests of the per-dataset correlations against zero, Holm-corrected over the seven diagnostics of each method set, and they treat the datasets as independent. Between the boundary arm and \RUS, on the other hand, the dataset-level G-mean loss tracks how much closer to the minority class the arm moves (Spearman correlation with the BP difference $+0.49$ for \MUSBS{} and $+0.51$ for \MUSBO, $p_\mathrm{H}\le0.001$) and not how much coverage it loses ($+0.04$ and $+0.02$). These contrasts suggest that coverage matters and that moving the selection towards the minority class is associated with the boundary arm's loss of accuracy, although the collinear diagnostics cannot establish either as a cause.

Better coverage, however, is not sufficient. The coverage arm covers the majority class better than \RUS{} on all 55 datasets and is closer to the full majority distribution on 51 (\MUSCS) and 48 (\MUSCO) of them, yet its dataset-level G-mean gain over \RUS{} is not associated with its coverage advantage (Spearman correlation with the CE difference $-0.08$ for \MUSCS{} and $-0.06$ for \MUSCO; Figure~\ref{fig:mechanism}d). Beyond the coverage that a random sample already achieves, more coverage is not associated with more accuracy; we refer to this as coverage saturation.

\paragraph{Threshold placement and ranking.} Part of the boundary arm's loss concerns where the scores fall relative to the default threshold. When the threshold is chosen on inner out-of-fold scores (Section~\ref{sec:design:revision}; first three repetitions, 15 folds per dataset), the G-mean deficit of \MUSBS{} to \RUS{} disappears for logistic regression (median $+0.000$, $p_\mathrm{H}=0.22$) and shrinks for SVM (from $-0.042$ to $-0.006$) and random forests (from $-0.017$ to $-0.008$), where it remains significant (Table~\ref{tab:x1}). It is unchanged for CART ($-0.023$), whose leaf frequencies leave the threshold almost no room to move (it moved in 1--4\% of the folds), and for $k$-NN ($-0.061$ before and $-0.067$ after adjustment). The sensitivity deficit falls accordingly, from $-0.142$ to $-0.013$ for SVM and from $-0.083$ to $-0.013$ for logistic regression. The ROC-AUC differences separate the two situations: for logistic regression and SVM the boundary arm ranks the test instances almost as well as \RUS{} (median differences $-0.0002$, not significant, and $-0.0012$), whereas for $k$-NN, CART and random forests it also ranks them worse ($-0.054$, $-0.021$ and $-0.018$). Tuning the hyperparameters at the default threshold reduces the deficit for $k$-NN, SVM and random forests, hardly changes it for CART and logistic regression, and removes it for no classifier. \MUSBO{} behaves in the same way (logistic regression after adjustment: $-0.003$, $p_\mathrm{H}=0.21$). \MUSCS{} is not significantly different from \RUS{} in G-mean in any of the three variants except tuned $k$-NN ($+0.003$, $p_\mathrm{H}=0.015$), and, averaged over the classifiers, it has the best average G-mean rank among the eight balancing methods in each variant (2.82, 2.92 and 2.92 for default, threshold-adjusted and tuned classifiers, against 3.10, 3.32 and 3.31 for \RUS{} and 3.05, 3.23 and 3.09 for ClusterCentroids).

\paragraph{Intervention.} Two variants of \MUSBS{} test the explanation offered above (Section~\ref{sec:design:revision}; Table~\ref{tab:x3}). Withholding majority partners from rare and outlier minority instances (safecap) raised sensitivity relative to \MUSBS{} only marginally (median at most $+0.002$, significant for all five classifiers) and lowered specificity, so that the mean G-mean change was negative for every classifier (mean $-0.015$ to $-0.035$, median $-0.003$ to $0.000$; significant for CART and random forests). On ten datasets with too few safe and borderline minority instances, this variant fell back to standard \MUSBS{} in 124 of the 1375 folds. Excluding majority instances with at least three minority instances among their five nearest neighbours from the matching (nounsafe) improved G-mean significantly for every classifier, but only slightly (mean $+0.006$ to $+0.010$); such majority instances were absent in half of the folds (mean 9.4 per fold). Both variants remained significantly worse than \RUS{} in G-mean and in sensitivity for every classifier. The partners of rare and outlier minority instances therefore do not account for the loss, and partners inside minority regions account for only a small part of it.

Taken together, the evidence points to two components of the boundary arm's loss. For logistic regression and SVM, and in part for random forests, the selection mainly shifts the scores relative to the default threshold, which threshold adjustment largely corrects. For $k$-NN, CART and random forests it also degrades the ranking of the test instances, which threshold adjustment cannot correct. Neither component is explained by the partners of atypical minority instances.

\input{tables/threshold_tuning}
\input{tables/intervention}

\input{tables/diag_summary}

\begin{figure}[t]
\centering
\includegraphics[width=\textwidth]{figures/fig3_mechanism.pdf}
\caption{Why the arms differ (55 datasets). (a) Boundary--coverage map of the nine balancing methods: median boundary proximity BP ($<1$: kept majority nearer the minority class) against median coverage error CE (logarithmic axis; lower = better coverage), with interquartile ranges; the stable and optimal variants coincide. NearMiss-3 keeps fewer majority instances than the target on 26 datasets. (b) Mean change in sensitivity and specificity relative to random undersampling, per classifier (small open marks; marker shape = classifier) and averaged over the classifiers (large marks, with 95\% bootstrap intervals over datasets); descriptive. (c) Within-dataset Spearman correlations across methods between G-mean (averaged over classifiers) and BP, between G-mean and CE, and between BP and CE: medians and interquartile ranges over the datasets, for the nine balancing methods and for the five boundary-focused methods (NearMiss-1 to 3, \MUSBS{} and \MUSBO). (d) Coverage saturation: dataset-level change in CE against change in G-mean (averaged over classifiers) of \MUSCS{} and \MUSCO{} relative to random undersampling; the grey band is $\pm0.01$. This figure also appears as Figure~4 of the main paper.}\label{fig:mechanism}
\end{figure}

\subsection{When does the boundary arm lose less?}\label{sec:results:when}

The loss of the boundary arm depends strongly on the composition of the minority class (Figure~\ref{fig:when}; Table~\ref{tab:meta_gain_corr}). It shrinks as the share of safe minority examples \citep{napierala2016types} grows (Spearman $\rho=+0.72$ for \MUSBS, Holm-adjusted $p<0.05$) and grows with the shares of outlier ($\rho=-0.68$) and rare ($\rho=-0.59$) minority examples and with class overlap (feature-overlap measure F1, $\rho=-0.61$, where larger values mean more overlap \citep{lorena2020how}; neighbourhood measure N3, $\rho=-0.46$). Fourteen datasets have no safe minority example; without them the correlation with the safe share is $+0.60$ for \MUSBS{} and $+0.58$ for \MUSBO{} ($n=41$, $p<0.001$), so the association does not rest on these datasets alone. The correlations are stronger on the 24 family representatives (safe share $+0.82$, outlier share $-0.78$), whereas size, dimensionality and imbalance ratio are not associated with the loss ($|\rho|\le0.14$). On the 19 datasets in which at least 80\% of the minority examples are safe, \MUSBS{} comes close to \RUS{} (mean difference $-0.013$, just outside the region of practical equivalence, median $-0.0001$; 11 datasets within $\pm0.01$, 6 below and 2 above); on the remaining 36 datasets it loses 0.090 on average. An exploratory depth-two regression tree (in-sample $R^2=0.63$) splits first on the outlier share (at 0.21) and then on the safe share (at 0.69) and on the share of majority instances in overlapping regions; its four leaves have mean differences of $-0.013$ (22 datasets with few outliers and mostly safe minority examples, a superset of the 19 datasets above), $-0.065$ (16 with few outliers but fewer safe examples), $-0.094$ (12 with many outliers but no overlapping majority instances) and $-0.203$ (5 with many outliers and overlap). Even in its favourable regime the boundary arm thus only approaches \RUS: it exceeds \RUS{} by more than 0.01 on 2 of the 55 datasets (by 0.018 to 0.019). For the coverage arm, whose gains over \RUS{} are centred at zero, no characteristic is associated with the gain over \RUS{} after Holm adjustment; its gain over the best non-matching method, however, increases with the share of rare minority examples (\MUSCS: $\rho=+0.48$, Holm-adjusted $p<0.05$; \MUSCO: $+0.32$, unadjusted $p<0.05$ only). The extreme datasets point in the same direction, although the correlations with the gain over \RUS{} are not significant: the largest classifier-averaged gains of the coverage arm over \RUS{} (up to $+0.036$) occur on datasets with almost no safe minority examples (\emph{abalone19}, three \emph{winequality} datasets and \emph{yeast\_1\_vs\_7}, safe share at most 0.07) and on \emph{shuttle\_6\_vs\_2\_3}, and the largest losses (down to $-0.038$) on two \emph{kr-vs-k} chess datasets, whose minority examples are all safe, and on two \emph{poker} datasets. The two \emph{kr-vs-k} datasets are the only ones on which the coverage arm differs from \RUS{} by more than two corrected standard errors (Table~\ref{tab:a6_per_dataset}), about the number expected by chance among 55 datasets.

\begin{figure}[t]
\centering
\includegraphics[width=\textwidth]{figures/fig4_when.pdf}
\caption{Dataset-level G-mean difference to random undersampling (averaged over the five classifiers) against the share of safe minority examples \citep{napierala2016types}, for the boundary arm (filled marks, with LOWESS smooths) and the coverage arm (open marks); one mark per dataset and method. The 14 datasets without safe minority examples are spread by a horizontal jitter of at most $\pm0.006$ for visibility (all fits use the unjittered values); the dashed line is the LOWESS smooth for \MUSBS{} fitted without them ($n=41$). The grey band is the region of practical equivalence ($\pm0.01$), and the dotted line marks a safe share of 0.8.}\label{fig:when}
\end{figure}

\subsection{Cost}\label{sec:results:runtime}

We timed every undersampler in isolation on the first training fold of every dataset, in one process with one thread and with five repetitions (Table~\ref{tab:runtime}; hardware in Section~\ref{sec:design:implementation}). The ordering of the methods agrees with the timings taken during the benchmark under parallel load, which were 1.5 to 3.5 times longer (Spearman correlation of the median times over the sixteen methods 0.997; per-dataset timings under load in Figure~\ref{fig:runtime_grid}). Figure~3 of the main paper sets the isolated median times against accuracy. \MUSBO{} (median 0.006\,s) and \MUSBS{} (0.010\,s) are as fast as NearMiss (0.010 to 0.012\,s). \MUSCS{} (0.14\,s) takes about four times as long as ClusterCentroids (0.033\,s) at the median, and its time grows faster with the majority size (log--log slopes 1.41 and 0.46), but it is faster than condensed nearest neighbour (0.71\,s) and the instance hardness threshold (2.56\,s). The exact \MUSCO{} is the expensive member of the family: its median time of 0.90\,s is 6.7 times that of \MUSCS, its time grows with a log--log slope of 2.15, and it reaches 506\,s on the largest training fold. On synthetic data with up to 100{,}000 majority instances (Table~\ref{tab:scaling}), \MUSBS{} needed 65\,s and \MUSCS{} 445\,s at the largest size, whereas \MUSBO{} exceeded the memory limit of its dense cost matrix and \MUSCO{} did not finish within 600\,s; at 30{,}000 majority instances the four methods took between 0.7 and 82\,s. Since \MUSCO{} brings no accuracy gain over \MUSCS, the stable (greedy) variant is the practical choice in the coverage arm.

\subsection{Robustness, secondary metrics and ablations}\label{sec:results:robustness}

The conclusions are robust to the choice of datasets (Table~\ref{tab:robustness}). On the 24 family representatives, the 46 datasets that are not separable by a single feature and the low- and high-IR groups, no significant difference changes direction; some lose significance in the smaller subsets. The only significant G-mean advantage of a matching undersampler appears for random forests on the 36 datasets with $\mathrm{IR}\ge9$, where \MUSCS{} ($+0.0027$, $p_\mathrm{H}=0.040$) and \MUSCO{} ($+0.0022$, $p_\mathrm{H}=0.043$) beat \RUS. Among the secondary metrics, the coverage arm is significantly better than \RUS{} in F1 and MCC for four of the five classifiers, and in PR-AUC for four (\MUSCO) and two (\MUSCS) classifiers. It is never significantly better than ClusterCentroids, \MUSCS{} is significantly worse than ClusterCentroids in PR-AUC for two classifiers, and the best non-matching methods (mostly cleaning methods or no resampling) are better in PR-AUC for every classifier. Compared with the non-undersampling references, the coverage arm has a higher G-mean than SMOTE and class weighting for CART and random forests and a lower one for logistic regression (Table~\ref{tab:reference_gmean}); in ROC-AUC it is better than class weighting for CART but worse than SMOTE for logistic regression and random forests and worse than class weighting for SVM, logistic regression and random forests (Table~\ref{tab:reference_auc}).

\paragraph{Row order and duplicates.} Rerunning the matching undersamplers on randomly permuted training rows (Section~\ref{sec:design:revision}) changed the boundary-arm selections on 27 (\MUSBS) and 26 (\MUSBO) datasets, 20 and 19 of which contain exact duplicate rows, and left them identical in about 70\% of the folds (median Jaccard index 1.0). The selections of the coverage arm overlapped only at chance level (median Jaccard index 0.036, against 0.032 expected by chance), because the seeded keep coin is attached to the pair members in row order, so that a permutation acts like a new seed; the round-one pairs were unchanged in the median fold. No permuted-versus-original difference was significant for any method, classifier or metric (smallest Holm-adjusted $p=0.15$), and all comparisons with \RUS{} kept their direction. With Holm adjustment over the three matching methods per classifier and metric, as in all additional experiments, \MUSCS{} was significantly better than \RUS{} in G-mean only for random forests in the original run ($+0.0016$) but for logistic regression ($+0.0006$) and random forests ($+0.0028$) after permutation, so the significance of this very small advantage depends on the realisation. Twenty-nine datasets contain exact duplicate rows (6622 rows in total, 122 of them in the minority class, at most 1235 in one dataset); four of them contain identical feature vectors with both labels, which we kept. After removing the duplicates and redrawing the folds (Table~\ref{tab:x6}), the boundary arm remained significantly worse than \RUS{} in G-mean for every classifier and in ROC-AUC for every classifier except logistic regression. The G-mean advantage of \MUSCS{} over \RUS{} became significant for SVM ($+0.0016$), logistic regression ($+0.0013$) and random forests ($+0.0027$) but stayed far inside the region of practical equivalence.


\paragraph{Dependence, aggregation and multiplicity.} Resampling the 24 base-data families instead of the 55 datasets widened the bootstrap standard errors of the mean paired differences by a median factor of 1.22 (range 0.64 to 1.51; Table~\ref{tab:a1_family_bootstrap}). All 20 family-clustered 95\% intervals of the boundary arm's G-mean difference to \RUS{} and ClusterCentroids lay below zero, whereas all mean differences of the coverage arm lay within 0.0075 of zero; the two bootstrap schemes disagreed only about intervals that touch zero, for example the boundary arm's ROC-AUC loss with logistic regression, whose family-clustered interval contains zero. Computing G-mean, F1 and MCC from confusion matrices pooled over the folds of each repetition, which avoids averaging metrics over test folds with two to four minority instances, changed none of the 120 test verdicts (Table~\ref{tab:a2_pooled}). In the primary family of 40 comparisons with \RUS, 25 were significant after Holm's correction: the boundary arm was worse in G-mean for all ten method--classifier pairs and in ROC-AUC for eight of them (not for logistic regression), whereas the coverage arm was never significantly different in G-mean and was better in ROC-AUC for seven of its ten pairs, with median gains of at most 0.0022.

The width of the region of practical equivalence matters for the coverage arm (Table~\ref{tab:a3_bayes_rope}). Its equivalence with \RUS{} in G-mean holds for all ten method--classifier pairs at $\pm0.01$ and $\pm0.02$, but for only two at $\pm0.005$ and for none at $\pm0.0025$, where the posteriors are mostly inconclusive; its equivalence with ClusterCentroids holds for six pairs at $\pm0.01$ and for all at $\pm0.02$. The benchmark therefore cannot exclude average G-mean differences of a quarter to half a percentage point between the coverage arm and \RUS. A width of $\pm0.01$ is close to the noise of a single dataset: the corrected resampled-$t$ standard error of the classifier-averaged G-mean difference between \MUSCS{} and \RUS{} has a median of 0.012 over the datasets (Table~\ref{tab:a6_per_dataset}). The boundary arm's G-mean deficit is decisive at every width, against both \RUS{} and ClusterCentroids, and it exceeds two standard errors on 25 datasets, 24 of them losses.

Raising the target majority:minority ratio from 1 to 1.5 or 2 increases the mean ROC-AUC in 17 of the 18 family--classifier combinations at each ratio (the exception is \MUSCS{} with CART, $-0.0001$ and $-0.0006$, not significant), significantly in 13, and never decreases it significantly (CART, $k$-NN and SVM; Table~\ref{tab:ablation} and Figure~\ref{fig:ablation}). The effect on G-mean depends on the classifier and the method: with SVM it falls significantly for every family except NearMiss-1, with $k$-NN it falls at ratio 2 for ClusterCentroids and \MUSBO, and with CART it rises slightly for the boundary arm at ratio 1.5 and markedly for NearMiss-1 at both ratios. The boundary arm is the most sensitive; for SVM, the G-mean of \MUSBS{} drops by 0.124 on average at ratio 1.5 ($p_\mathrm{H}=6.1\times10^{-6}$). Keeping the pair member closer to the minority class in \MUSCS{} and using class-conditional Mahalanobis preferences in \MUSBS{} never improve significantly on the defaults in G-mean or ROC-AUC: the boundary keep rule changes G-mean for none of the three classifiers and lowers ROC-AUC for one, and the asymmetric preferences lower G-mean for one classifier and ROC-AUC for two. The boundary keep rule does, however, improve F1 and MCC for all three classifiers and PR-AUC for two: it raises precision and specificity at the cost of sensitivity (median precision differences $+0.009$ to $+0.021$, sensitivity differences $-0.003$ to $-0.025$). Because the coverage methods retain nearly identical distributions (Section~\ref{sec:results:factorial}), we expect the \MUSCS{} ablations to carry over to \MUSCO, whose ablations we did not run because of its cost.

\subsection{Spatially balanced sampling, ensembles and seed variability}\label{sec:results:lpm}

The coverage arm spreads the retained majority instances over the majority class, as spatially balanced sampling designs do. The local pivotal method \citep{grafstrom2012spatially} with equal inclusion probabilities performed like \MUSCS: no difference between the two was significant for any classifier in G-mean or ROC-AUC (largest absolute median difference 0.0003; Table~\ref{tab:x4}). Like \MUSCS, it was marginally better than \RUS{} (significant in ROC-AUC for $k$-NN, SVM, logistic regression and random forests and in G-mean for logistic regression; largest median difference $+0.0023$), whereas random undersampling stratified by $k$-means clusters could not be distinguished from \RUS. \MUSCS{} had the lowest coverage error (median 1.56, against 1.61 for the local pivotal method, 1.74 for stratified and 1.81 for plain random undersampling; the local pivotal method was lower than \MUSCS{} on only 8 of the 55 datasets), and the local pivotal method had the best spatial balance (median Voronoi index 0.23, against 0.29 for \MUSCS, 0.32 for stratified and 0.46 for plain random undersampling). In performance and in coverage, the coverage arm thus behaves like an established spatially balanced design.

Averaging the scores of ten classifiers trained on different undersampled subsets improved ROC-AUC over a single subset for all three samplers and every classifier, most for CART (G-mean about $+0.02$, ROC-AUC about $+0.06$; Table~\ref{tab:x5}), and improved G-mean for every combination except \MUSCS{} and the local pivotal method with $k$-NN. Ensembles of \MUSCS{} or local-pivotal subsets were not better than ensembles of random subsets; the only significant difference favoured \RUS{} ($k$-NN, G-mean $-0.002$). The ten members also show how much a single subset depends on its seed. The across-seed standard deviation of G-mean was significantly lower for \MUSCS{} and the local pivotal method than for \RUS{} with $k$-NN, SVM, logistic regression and random forests (medians 0.019, 0.011, 0.011 and 0.018 for \MUSCS{} against 0.023, 0.015, 0.017 and 0.023 for \RUS, i.e.\ 20 to 35\% lower), but not with CART (0.036 against 0.039); the same holds for ROC-AUC, where the reduction is also significant for CART. This lower variability of a single subset is the most tangible practical difference between the coverage arm and random undersampling, and it disappears once several subsets are averaged.

\input{tables/spatial_ensembles}
